\documentclass[11pt,a4paper]{article}
\usepackage[margin=20mm]{geometry}
\usepackage[T1]{fontenc}
\usepackage{lmodern,amsmath,amssymb,graphicx,booktabs,microtype}
\usepackage[hidelinks]{hyperref}
\input{glyphtounicode}
\pdfgentounicode=1
\hypersetup{pdftitle={Inferring a ball's launch angle with a normalizing flow}}
\setlength{\parindent}{0pt}
\setlength{\parskip}{4pt}
\setlength{\abovedisplayskip}{7pt}
\setlength{\belowdisplayskip}{7pt}
\raggedbottom
\newcommand{\E}{\mathbb E}
\newcommand{\Var}{\operatorname{Var}}
\newcommand{\KL}{D_{\mathrm{KL}}}
\newcommand{\sigmoid}{\operatorname{sigmoid}}
\newcommand{\Ang}{\mathsf A}
\newcommand{\figpath}{results/figures/}
\newcommand{\BestEpoch}{31}
\newcommand{\EpochsRun}{46}
\newcommand{\TrainSeconds}{39.3}
\newcommand{\BestValidation}{2.382}
\newcommand{\MaxMeanChange}{4.73e-10}
\newcommand{\MaxBranchChange}{7.09e-12}
\newcommand{\MaxTVChange}{1.92e-05}
\newcommand{\MaxMassChange}{4.23e-05}


\begin{document}
\begin{center}
{\LARGE\bfseries Inferring a ball's launch angle}\\[5pt]
{\large A worked example of neural posterior estimation with a normalizing flow}\\[4pt]
{\small STA492: simulation, probability, derivations, and measured results}
\end{center}

\section*{1. From an angle to a landing distance}
We observe one number, the distance at which a ball lands, and want a
distribution over its launch angle. Different angles can give similar ranges;
the answer can therefore have two peaks. This example adapts the introductory
ball simulator in Deistler et al.\ \cite{guide} and accompanies Weniger's SBI
lecture \cite{lecture}. We omit air resistance to make the physics explicit,
use SI constants, and choose Gaussian measurement noise. These choices differ
from the guide's numerical code \cite{guidecode}; this is not an exact reproduction.

Write $\alpha$ for the angle in degrees and $\theta=\pi\alpha/180$ for the
same angle in radians. The initial speed is $v_0=13\,\mathrm{m/s}$ and gravity
is $g=9.81\,\mathrm{m/s^2}$. During one throw, $w\geq0$ is a constant
horizontal acceleration caused by tailwind. With no drag,
\[
 \ddot X(t)=w,\qquad \ddot Y(t)=-g,
 \quad X(0)=Y(0)=0,\quad
 \dot X(0)=v_0\cos\theta,\quad \dot Y(0)=v_0\sin\theta.
\]
Integrating once gives $\dot X(t)=v_0\cos\theta+wt$ and
$\dot Y(t)=v_0\sin\theta-gt$. Integrating again gives
\[
 X(t)=v_0\cos\theta\,t+\tfrac12wt^2,\qquad
 Y(t)=v_0\sin\theta\,t-\tfrac12gt^2.
\]
To find the landing time, factor the height equation:
$Y(T)=T(v_0\sin\theta-gT/2)=0$. The root $T=0$ is the start; the other is
\[
 T(\alpha)=\frac{2v_0\sin\theta}{g}.
\]
Substitute it into $X(t)$ and use $2\sin\theta\cos\theta=\sin(2\theta)$:
\begin{equation}
 \boxed{D(\alpha,w)=A(\alpha)+B(\alpha)w,\quad
 A(\alpha)=\frac{v_0^2}{g}\sin(2\theta),\quad
 B(\alpha)=\tfrac12T(\alpha)^2.}
 \label{eq:range}
\end{equation}
$A$ is a distance, $B$ has units $\mathrm{s^2}$, and $Bw$ is a distance.
Wind here is an acceleration, not a wind speed. The simulator calculates this
formula directly; it does not discretize time. Trajectories illustrate the
physics, but only the final measured distance is used for inference.

\begin{center}
\includegraphics[width=0.97\linewidth]{\figpath trajectories.pdf}\\[-2pt]
{\small Two no-wind angles can share a range; different winds change the range at a fixed angle.}
\end{center}

\newpage
\section*{2. From deterministic physics to a probability model}
Before observing a distance, use a Normal angle distribution truncated to
$0<\alpha<90$ degrees. Let $\varphi$ and $\Phi$ be the standard Normal density
and CDF. Its normalized density per degree is
\[
 p(\alpha)=\frac{\varphi((\alpha-35)/15)}{15C}\,\mathbf1_{(0,90)}(\alpha),
 \qquad C=\Phi\!\left(\frac{90-35}{15}\right)-\Phi\!\left(\frac{-35}{15}\right).
\]
The numbers 35 and 15 describe the Normal before truncation. Each simulated
throw draws independent wind and measurement error:
\[
 W\sim\operatorname{Exp}(\lambda),\quad \lambda=2,
 \qquad \varepsilon\sim N(0,\sigma_x^2),\quad \sigma_x=0.25\,\mathrm m,
 \qquad x=A(\alpha)+B(\alpha)W+\varepsilon.
\]
Thus $\E W=1/\lambda=0.5\,\mathrm{m/s^2}$ and
$\Var W=1/\lambda^2$. Independence and \eqref{eq:range} give useful checks:
\[
 \E[x\mid\alpha]=A(\alpha)+B(\alpha)/\lambda,\qquad
 \Var(x\mid\alpha)=\sigma_x^2+B(\alpha)^2/\lambda^2.
\]
Given angle and wind, the density is Gaussian:
$p(x\mid\alpha,w)=N(x;A(\alpha)+B(\alpha)w,\sigma_x^2)$.
We infer only the angle, so the likelihood must average over unobserved wind:
\begin{equation}
 p(x\mid\alpha)=\int_0^\infty
 N(x;A(\alpha)+B(\alpha)w,\sigma_x^2)\,\lambda e^{-\lambda w}\,dw.
 \label{eq:marginal}
\end{equation}
Sampling wind inside the simulator generates this marginal distribution
automatically. Holding it fixed would define a different inference problem.
Bayes' theorem gives
\begin{equation}
 p(\alpha\mid x_o)=\frac{p(x_o\mid\alpha)p(\alpha)}{p(x_o)},\qquad
 p(x_o)=\int_0^{90}p(x_o\mid a)p(a)\,da.
 \label{eq:bayes}
\end{equation}
Without wind or measurement noise, solve $x_o=(v_0^2/g)\sin(2\theta)$.
For $0<gx_o/v_0^2<1$, the two solutions are
\[
 \theta_1=\tfrac12\arcsin(gx_o/v_0^2),\qquad
 \theta_2=\pi/2-\theta_1.
\]
For $x_o=13\,\mathrm m$ these are $24.50^\circ$ and $65.50^\circ$.
Noise broadens the possibilities. The asymmetric prior and angle-dependent
wind effect change the branches' positions and weights; the noisy posterior
is not simply two equal spikes at complementary angles.

\begin{center}
\includegraphics[width=0.94\linewidth]{\figpath prior_and_joint.pdf}\\[-2pt]
{\small Prior angles generate joint pairs $(\alpha,x)$; observations select conditional angle distributions.}
\end{center}

\newpage
\section*{3. Why learning from simulations targets a posterior}
An NPE dataset is generated in the forward direction:
$\alpha_i\sim p(\alpha)$, followed by $x_i\sim p(x\mid\alpha_i)$.
Each pair follows $p(\alpha,x)=p(\alpha)p(x\mid\alpha)$.
The neural network represents a normalized conditional density
$q_\phi(\alpha\mid x)$, with trainable weights $\phi$. Minimize
\[
 \widehat{\mathcal L}(\phi)=-\frac1N\sum_{i=1}^N\log q_\phi(\alpha_i\mid x_i).
\]
The network receives a distance and is rewarded for assigning higher density
to the angle that generated it. For continuous variables this is a density,
not a probability of equality to one exact angle. The target is a distribution,
not a mean-squared-error prediction of a single angle.

Let $\Ang$ denote the random angle in degrees and $X$ the random distance.
Using $p(\alpha,x)=p(x)p(\alpha\mid x)$, the population loss is
\begin{align*}
 \mathcal L(\phi)
 &=-\int p(x)\int p(\alpha\mid x)\log q_\phi(\alpha\mid x)\,d\alpha\,dx\\
 &=\int p(x)\int p(\alpha\mid x)
 \left[-\log p(\alpha\mid x)
 +\log\frac{p(\alpha\mid x)}{q_\phi(\alpha\mid x)}\right]d\alpha\,dx\\
 &=\underbrace{H_p(\Ang\mid X)}_{\text{independent of }\phi}
 +\E_{p(x)}\!\left[\KL\{p(\alpha\mid x)\Vert q_\phi(\alpha\mid x)\}\right].
\end{align*}
The second line adds and subtracts the true posterior log-density. KL is
nonnegative: $\log t\leq t-1$ implies
$\int p\log(p/q)\geq\int(p-q)=0$ on the common support.
Hence minimizing the population NPE loss minimizes average posterior KL.
Equality with the true posterior requires sufficient model capacity and an
attained global population optimum, for almost every prior-predictive $x$.
Finite simulations, limited capacity, and imperfect optimization leave error.

\textbf{The actual training procedure.}
We generate 20,000 pairs once (NumPy seed 2026), split them randomly into
18,000 training and 2,000 validation pairs, and standardize distances using
training statistics only. With PyTorch seed 42, Adam uses learning rate
$10^{-3}$, batches of 256, and gradient norm clipping at 5. There are at most
100 epochs; an epoch is one pass through the training pairs. We stop after
15 epochs without a validation improvement of $10^{-4}$ and restore the
best validation weights. No reference likelihood is used during training.
The supplied run completed \EpochsRun{} epochs, restored epoch \BestEpoch{},
and took \TrainSeconds{} CPU seconds with two PyTorch threads.

\begin{center}
\includegraphics[width=0.80\linewidth]{\figpath training.pdf}\\[-2pt]
{\small Training optimizes simulated-pair log density; validation selects the saved weights.}
\end{center}
The same trained network then handles 13, 15, and 16 metres without further
optimization. This reuse is called \emph{amortized inference}.

\newpage
\section*{4. How a normalizing flow represents two angle peaks}
A flow generates samples by transforming $z\sim N(0,1)$ through an invertible
map. Here $z\to u$ consists of three increasing rational-quadratic spline
maps conditioned on the distance \cite{spline,nflows}. Finally,
\[
 \alpha=90\sigmoid(u),\qquad
 u=\log\frac{\alpha}{90-\alpha},\qquad
 \frac{d\alpha}{du}=90\sigmoid(u)\sigmoid(-u).
\]
This enforces the angle bounds without clipping. Probability is conserved:
$q_\phi(\alpha\mid x)\,d\alpha=p_Z(z)\,dz$. Therefore
\begin{equation}
 \log q_\phi(\alpha\mid x)=\log p_Z(z)
 +\log\left|\frac{dz}{d\alpha}\right|.
 \label{eq:flow}
\end{equation}
Density evaluation runs angle $\to u\to z$, the inverse of sampling.
In the code, spline \texttt{forward} means the density-evaluation direction,
and \texttt{inverse=True} means the sampling direction. With
$f=\alpha/90$, the angle-to-$u$ log-Jacobian is
$-\log90-\log f-\log(1-f)$. We add each spline log-Jacobian along the chain.
All densities are per degree; changing angle units would require its own
Jacobian. We assign density zero outside the box and at the endpoints,
which have probability zero.

\textbf{One spline bin, explicitly.}
Consider an increasing map $S$ taking $u_k$ to $v_k$ and $u_{k+1}$ to
$v_{k+1}$. Define bin width $w_k=u_{k+1}-u_k>0$, height
$h_k=v_{k+1}-v_k>0$, secant slope $\delta_k=h_k/w_k$, and
$\xi=(u-u_k)/w_k\in[0,1]$. Let $d_k,d_{k+1}>0$ be endpoint derivatives.
The bin formula is
\begin{equation}
 S(u)=v_k+h_k\,
 \frac{\delta_k\xi^2+d_k\xi(1-\xi)}
 {\delta_k+(d_{k+1}+d_k-2\delta_k)\xi(1-\xi)}.
 \label{eq:spline}
\end{equation}
At $\xi=0$ it equals $v_k$; at $\xi=1$ it equals $v_k+h_k=v_{k+1}$.
Apply the quotient rule and $d\xi/du=1/w_k$ to obtain
\begin{equation}
 S'(u)=\frac{\delta_k^2
 [d_{k+1}\xi^2+2\delta_k\xi(1-\xi)+d_k(1-\xi)^2]}
 {[\delta_k+(d_{k+1}+d_k-2\delta_k)\xi(1-\xi)]^2}>0.
 \label{eq:derivative}
\end{equation}
This positivity ensures invertibility. The log absolute derivative in
\eqref{eq:flow} is cheap to evaluate. Inversion solves the quadratic obtained
by rearranging \eqref{eq:spline}; \texttt{nflows} handles the root calculation.

Each layer has 16 bins. Its network maps one standardized distance through
two 32-unit tanh hidden layers to 47 outputs: 16 raw widths, 16 raw heights,
and 15 raw interior derivatives. Softmax makes widths/heights positive;
softplus plus a small positive floor makes derivatives positive. The spline
maps $[-5,5]$ onto itself and is the identity outside; endpoint derivatives
are fixed to one. Initialization uses equal widths/heights and unit slopes.
Only this spline arithmetic is borrowed from the library; the conditioning
networks, bounded-angle transformation, training, and sampling are explicit.

An increasing map need not yield a unimodal density. In the sampling
direction, $q_\phi(\alpha\mid x)=p_Z(z)/|d\alpha/dz|$: different regions
can squeeze samples more tightly, creating two density peaks while keeping
their order. Normalization of the model alone does not imply posterior accuracy.

\begin{center}
\includegraphics[width=0.90\linewidth]{\figpath flow_map.pdf}\\[-2pt]
{\small The sampling map changes sample spacing, producing a density with two angle peaks.}
\end{center}

\newpage
\section*{5. What the trained network learned}
The three distances are hypothetical observations, with no single known
generating angle. A numerical reference evaluates \eqref{eq:marginal} using
adaptive quadrature. Substituting $q=\lambda w$ gives the convenient integrand
$e^{-q}N(x_o;A+Bq/\lambda,\sigma_x^2)$ over $q\in[0,\infty)$.
Multiplying by the prior and normalizing over the full angle interval gives
\eqref{eq:bayes}. This calculation checks the trained model; it never supplies
training targets or gradients.

\begin{center}
\includegraphics[width=\linewidth]{\figpath posterior_comparison.pdf}\\[-2pt]
{\small One trained flow (dashed blue), numerical references (black), and the prior (gray).}
\end{center}

Define $P_{<45}=P(\alpha<45^\circ\mid x_o)$. This is a fixed low-angle
partition, not a fitted boundary between modes. Total variation is
\[
 \operatorname{TV}(p,q)=\tfrac12\int_0^{90}|p(\alpha\mid x_o)-q_\phi(\alpha\mid x_o)|\,d\alpha.
\]
It ranges from zero to one. For midpoint grid spacing $\Delta\alpha$, we
approximate it by $\tfrac12\sum_j|p_j-q_j|\Delta\alpha$. The learned density
is not renormalized on the grid; its original integral is reported separately.

\begin{center}
\small
\begin{tabular}{@{}rrrrrr@{}}
\toprule
$x_o$ (m) & Ref.\ modes ($^\circ$) & Ref.\ $P_{<45}$ & Flow $P_{<45}$ & TV & Flow integral\\
\midrule
\input results/accuracy_rows.tex
\bottomrule
\end{tabular}\\[3pt]
{\small Densities use a 4,000-point grid; posterior sampling uses 20,000 draws per observation.}
\end{center}

Reference grids of 2,000 and 4,000 points change the mean by at most
\texttt{\MaxMeanChange} degrees and $P_{<45}$ by at most
\texttt{\MaxBranchChange}. Learned-grid refinement changes TV by at most
\texttt{\MaxTVChange} and the integral by at most \texttt{\MaxMassChange}.
An integral near one checks quadrature, not agreement with the reference.
The complete values and wind integration error estimates are saved with the run.

All three learned densities retain two peaks, with small location shifts and
slightly too much probability below $45^\circ$. Thus the qualitative ambiguity
is recovered, but the learned answer is still approximate.
The prior favors angles near $35^\circ$, so low and high branches need not
receive equal weight. A posterior mean can fall in the valley between peaks,
and a central credible interval can include that valley. Both can obscure
which angles actually have high posterior density. Compare the full density
and branch weights, including visible shifts or missing mass, before interpreting
one summary as the inferred launch.

\newpage
\section*{6. Predictive checks, verification, and reproducibility}
For a future throw with the same unknown angle but independent wind and noise,
the posterior predictive distribution is
\[
 p(x_{\rm new}\mid x_o)=\int p(x_{\rm new}\mid\alpha)
 p(\alpha\mid x_o)\,d\alpha.
\]
We approximate it by drawing $\alpha$ from the flow posterior and then calling
the simulator with a \emph{new} wind and measurement error. We do not reuse
the latent wind that could have produced the original observation.

\begin{center}
\includegraphics[width=0.95\linewidth]{\figpath posterior_predictive.pdf}\\[-2pt]
{\small Posterior predictive landing distances; dashed lines mark the conditioning observations.}
\end{center}

A future throw need not reproduce $x_o$ exactly. Predictive width comes from
uncertain angles and new random effects. These plots check consistency under
our assumed model, not agreement with independent real observations. A good
predictive plot alone cannot certify a correct angle posterior, since different
angles can predict similar ranges.

\textbf{What was tested.}
Ten tests check the landing time and no-wind range, complementary-angle
symmetry, the integrated wind effect, simulated conditional means/variances,
prior support and normalization, reference refinement, flow inversion,
Jacobian agreement with automatic differentiation, density integration,
finite training gradients, and saved standardization. The executed notebook
also checks support, learned integrals, grid refinement, and density agreement
after reloading the checkpoint.

\textbf{Limits.}
This is one training run and three distances, not a repeated-dataset coverage
study. Exact change of variables does not remove neural approximation error.
Air resistance, time-varying wind, uncertain speed, and uncertain gravity are
omitted. Gaussian noise can yield negative recorded ranges; we do not clip them.
The posterior answers the assumed model's question. New data under that model
can reuse the network; changing the simulator or prior requires retraining.

\textbf{How to reproduce.}
Start with \texttt{projectile\_demo.ipynb}; the README gives setup and run commands.
Use \texttt{python3 run\_notebook.py} for a fresh kernel and
\texttt{python3 -m unittest -v} for tests. Settings, seeds, versions, weights,
arrays, metrics, and figures are in \texttt{results/}. The notebook generates the
report table; rebuild the PDF and review its interpretation after a changed run.

\begingroup
\small\raggedright
\begin{thebibliography}{9}
\bibitem{guide} M. Deistler et al.\ (2025). \emph{Simulation-Based Inference: A Practical Guide}.
\href{https://arxiv.org/abs/2508.12939}{arXiv:2508.12939}, Section 2 and Figure 2.
\bibitem{guidecode} SBI developers. \emph{Ball-throw example accompanying the practical guide}.
\href{https://github.com/sbi-dev/sbi-practical-guide/blob/main/examples/ball_throw.ipynb}{Source notebook}.
\bibitem{lecture} C. Weniger (2026). \emph{Lecture 1: ABC and Neural SBI}. IAIFI PhD School.
\href{https://cweniger.github.io/teaching-2608-IAIFI/lecture1.html}{Teaching material}.
\bibitem{spline} C. Durkan, A. Bekasov, I. Murray, and G. Papamakarios (2019).
\emph{Neural Spline Flows}. NeurIPS. \href{https://arxiv.org/abs/1906.04032}{arXiv:1906.04032}.
\bibitem{nflows} C. Durkan et al.\ (2020). \emph{nflows: normalizing flows in PyTorch}, v0.14.
\href{https://github.com/bayesiains/nflows}{Software repository}.
\end{thebibliography}
\endgroup
\end{document}
